% fig14_cch_access_structure.tex — 图14 CCH 双通道访问结构与三堵墙
% 《深度学习与大语言模型：前沿技术全景报告》图示库
% 由 dl_llm_report.tex 抽取，经 % fig14_cch_access_structure.tex — 图14 CCH 双通道访问结构与三堵墙
% 《深度学习与大语言模型：前沿技术全景报告》图示库
% 由 dl_llm_report.tex 抽取，经 % fig14_cch_access_structure.tex — 图14 CCH 双通道访问结构与三堵墙
% 《深度学习与大语言模型：前沿技术全景报告》图示库
% 由 dl_llm_report.tex 抽取，经 % fig14_cch_access_structure.tex — 图14 CCH 双通道访问结构与三堵墙
% 《深度学习与大语言模型：前沿技术全景报告》图示库
% 由 dl_llm_report.tex 抽取，经 \input{figs/fig14_cch_access_structure} 引用（caption/label 在主文件）
% 注意：本文件为片段，依赖主文件导言区的 tikz/pgfplots 宏包、颜色定义与 tikzset 样式，不可单独编译。

\begin{tikzpicture}[x=1cm, y=1cm]
  % ---- 左段：信息源与瓶颈 ----
  \node[blk, align=center] (stream) at (0,0)
    {\textbf{无限 token 流}\\ $\cdots\bullet$ 苹果落地 $\bullet\cdots$\\ {\scriptsize 查询：跟踪 $+$ 回查}};
  \node[orn, align=center] (sigma) at (3.6,0)
    {\textbf{可访问状态} $\Sigma$\\ 瓶颈容量 $B$ 比特\\ {\scriptsize 固定每 token 预算}};
  % ---- 中段：双通道（上下分层） ----
  \node[grn, align=center] (cha) at (7.7,1.3)
    {\textbf{压缩态通道}\\ $O(1)$：SSM/RNN 递归态\\ {\scriptsize 擅长跟踪，弱逐字回查}};
  \node[dgt, align=center] (chb) at (7.7,-1.3)
    {\textbf{逐字索引通道}\\ KV Cache · 外置索引\\ {\scriptsize 擅长回查，弱窗外复合}};
  % ---- 三堵墙：间距 1.15cm，名称水平标于墙顶（虚线不穿任何文字） ----
  \foreach \x/\nm in {10.35/香农墙, 11.5/视界墙, 12.65/电路墙}{
    \draw[dashed, very thick, red!65!black] (\x,-2.35) -- (\x,2.35);
    \node[font=\scriptsize\bfseries, text=red!65!black] at (\x,2.62) {\nm};
  }
  % ---- 右段：完备混合体 ----
  \node[grn, align=center] (hybrid) at (14.9,0)
    {\textbf{访问完备混合体}\\ 双通道，按墙付费\\ $\Rightarrow$ 能力超可加};
  % ---- 连线：双通道斜穿三墙（对应“按墙付费、逐一穿越”） ----
  \draw[brr] (stream) -- (sigma);
  \draw[brr] (sigma.east) -- (cha.west);
  \draw[brr] (sigma.east) -- (chb.west);
  \draw[brr] (cha.east) -- ([yshift=2.5mm]hybrid.west);
  \draw[brr] (chb.east) -- ([yshift=-2.5mm]hybrid.west);
\end{tikzpicture}
 引用（caption/label 在主文件）
% 注意：本文件为片段，依赖主文件导言区的 tikz/pgfplots 宏包、颜色定义与 tikzset 样式，不可单独编译。

\begin{tikzpicture}[x=1cm, y=1cm]
  % ---- 左段：信息源与瓶颈 ----
  \node[blk, align=center] (stream) at (0,0)
    {\textbf{无限 token 流}\\ $\cdots\bullet$ 苹果落地 $\bullet\cdots$\\ {\scriptsize 查询：跟踪 $+$ 回查}};
  \node[orn, align=center] (sigma) at (3.6,0)
    {\textbf{可访问状态} $\Sigma$\\ 瓶颈容量 $B$ 比特\\ {\scriptsize 固定每 token 预算}};
  % ---- 中段：双通道（上下分层） ----
  \node[grn, align=center] (cha) at (7.7,1.3)
    {\textbf{压缩态通道}\\ $O(1)$：SSM/RNN 递归态\\ {\scriptsize 擅长跟踪，弱逐字回查}};
  \node[dgt, align=center] (chb) at (7.7,-1.3)
    {\textbf{逐字索引通道}\\ KV Cache · 外置索引\\ {\scriptsize 擅长回查，弱窗外复合}};
  % ---- 三堵墙：间距 1.15cm，名称水平标于墙顶（虚线不穿任何文字） ----
  \foreach \x/\nm in {10.35/香农墙, 11.5/视界墙, 12.65/电路墙}{
    \draw[dashed, very thick, red!65!black] (\x,-2.35) -- (\x,2.35);
    \node[font=\scriptsize\bfseries, text=red!65!black] at (\x,2.62) {\nm};
  }
  % ---- 右段：完备混合体 ----
  \node[grn, align=center] (hybrid) at (14.9,0)
    {\textbf{访问完备混合体}\\ 双通道，按墙付费\\ $\Rightarrow$ 能力超可加};
  % ---- 连线：双通道斜穿三墙（对应“按墙付费、逐一穿越”） ----
  \draw[brr] (stream) -- (sigma);
  \draw[brr] (sigma.east) -- (cha.west);
  \draw[brr] (sigma.east) -- (chb.west);
  \draw[brr] (cha.east) -- ([yshift=2.5mm]hybrid.west);
  \draw[brr] (chb.east) -- ([yshift=-2.5mm]hybrid.west);
\end{tikzpicture}
 引用（caption/label 在主文件）
% 注意：本文件为片段，依赖主文件导言区的 tikz/pgfplots 宏包、颜色定义与 tikzset 样式，不可单独编译。

\begin{tikzpicture}[x=1cm, y=1cm]
  % ---- 左段：信息源与瓶颈 ----
  \node[blk, align=center] (stream) at (0,0)
    {\textbf{无限 token 流}\\ $\cdots\bullet$ 苹果落地 $\bullet\cdots$\\ {\scriptsize 查询：跟踪 $+$ 回查}};
  \node[orn, align=center] (sigma) at (3.6,0)
    {\textbf{可访问状态} $\Sigma$\\ 瓶颈容量 $B$ 比特\\ {\scriptsize 固定每 token 预算}};
  % ---- 中段：双通道（上下分层） ----
  \node[grn, align=center] (cha) at (7.7,1.3)
    {\textbf{压缩态通道}\\ $O(1)$：SSM/RNN 递归态\\ {\scriptsize 擅长跟踪，弱逐字回查}};
  \node[dgt, align=center] (chb) at (7.7,-1.3)
    {\textbf{逐字索引通道}\\ KV Cache · 外置索引\\ {\scriptsize 擅长回查，弱窗外复合}};
  % ---- 三堵墙：间距 1.15cm，名称水平标于墙顶（虚线不穿任何文字） ----
  \foreach \x/\nm in {10.35/香农墙, 11.5/视界墙, 12.65/电路墙}{
    \draw[dashed, very thick, red!65!black] (\x,-2.35) -- (\x,2.35);
    \node[font=\scriptsize\bfseries, text=red!65!black] at (\x,2.62) {\nm};
  }
  % ---- 右段：完备混合体 ----
  \node[grn, align=center] (hybrid) at (14.9,0)
    {\textbf{访问完备混合体}\\ 双通道，按墙付费\\ $\Rightarrow$ 能力超可加};
  % ---- 连线：双通道斜穿三墙（对应“按墙付费、逐一穿越”） ----
  \draw[brr] (stream) -- (sigma);
  \draw[brr] (sigma.east) -- (cha.west);
  \draw[brr] (sigma.east) -- (chb.west);
  \draw[brr] (cha.east) -- ([yshift=2.5mm]hybrid.west);
  \draw[brr] (chb.east) -- ([yshift=-2.5mm]hybrid.west);
\end{tikzpicture}
 引用（caption/label 在主文件）
% 注意：本文件为片段，依赖主文件导言区的 tikz/pgfplots 宏包、颜色定义与 tikzset 样式，不可单独编译。

\begin{tikzpicture}[x=1cm, y=1cm]
  % ---- 左段：信息源与瓶颈 ----
  \node[blk, align=center] (stream) at (0,0)
    {\textbf{无限 token 流}\\ $\cdots\bullet$ 苹果落地 $\bullet\cdots$\\ {\scriptsize 查询：跟踪 $+$ 回查}};
  \node[orn, align=center] (sigma) at (3.6,0)
    {\textbf{可访问状态} $\Sigma$\\ 瓶颈容量 $B$ 比特\\ {\scriptsize 固定每 token 预算}};
  % ---- 中段：双通道（上下分层） ----
  \node[grn, align=center] (cha) at (7.7,1.3)
    {\textbf{压缩态通道}\\ $O(1)$：SSM/RNN 递归态\\ {\scriptsize 擅长跟踪，弱逐字回查}};
  \node[dgt, align=center] (chb) at (7.7,-1.3)
    {\textbf{逐字索引通道}\\ KV Cache · 外置索引\\ {\scriptsize 擅长回查，弱窗外复合}};
  % ---- 三堵墙：间距 1.15cm，名称水平标于墙顶（虚线不穿任何文字） ----
  \foreach \x/\nm in {10.35/香农墙, 11.5/视界墙, 12.65/电路墙}{
    \draw[dashed, very thick, red!65!black] (\x,-2.35) -- (\x,2.35);
    \node[font=\scriptsize\bfseries, text=red!65!black] at (\x,2.62) {\nm};
  }
  % ---- 右段：完备混合体 ----
  \node[grn, align=center] (hybrid) at (14.9,0)
    {\textbf{访问完备混合体}\\ 双通道，按墙付费\\ $\Rightarrow$ 能力超可加};
  % ---- 连线：双通道斜穿三墙（对应“按墙付费、逐一穿越”） ----
  \draw[brr] (stream) -- (sigma);
  \draw[brr] (sigma.east) -- (cha.west);
  \draw[brr] (sigma.east) -- (chb.west);
  \draw[brr] (cha.east) -- ([yshift=2.5mm]hybrid.west);
  \draw[brr] (chb.east) -- ([yshift=-2.5mm]hybrid.west);
\end{tikzpicture}
